\documentclass[12pt]{article}
\usepackage[T1]{fontenc}
\usepackage[utf8]{inputenc}
\usepackage{lmodern}
\usepackage{textcomp}
\usepackage{enumitem}
\usepackage{geometry}
\usepackage{graphicx}
\usepackage{amsmath, amssymb}
\geometry{margin=1in}
\begin{document}
\begin{center}
Physics - Undergraduate Level Assessment 21 \
Total Marks: 40 \
Answer all questions.
\end{center}

\section*{Multiple Choice (10 marks)}
\begin{enumerate}[label=\arabic*.]
\item A satellite of mass m orbits a planet of mass M in a circular orbit of radius R. The time required for one revolution is
\begin{enumerate}[label=(\alph*)]
    \item independent of M
    \item proportional to $m^(1$/2)
    \item linear in R
    \item proportional to $R^(3$/2)
\end{enumerate}
\item The energy required to remove both electrons from the helium atom in its ground state is 79.0 eV. How much energy is required to ionize helium (i.e., to remove one electron)?
\begin{enumerate}[label=(\alph*)]
    \item 24.6 eV
    \item 39.5 eV
    \item 51.8 eV
    \item 54.4 eV
\end{enumerate}
\item A grating spectrometer can just barely resolve two wavelengths of 500 nm and 502 nm, respectively. Which of the following gives the resolving power of the spectrometer?
\begin{enumerate}[label=(\alph*)]
    \item 2
    \item 250
    \item 5,000
    \item 10,000
\end{enumerate}
\item The eigenvalues of a Hermitian operator are always
\begin{enumerate}[label=(\alph*)]
    \item real
    \item imaginary
    \item degenerate
    \item linear
\end{enumerate}
\item The negative muon, mu^-, has properties most similar to which of the following?
\begin{enumerate}[label=(\alph*)]
    \item Electron
    \item Meson
    \item Photon
    \item Boson
\end{enumerate}
\end{enumerate}

\section*{True / False (10 marks)}
\textit{Answer True or False}
\begin{enumerate}[label=\arabic*.]
\setcounter{enumi}{5}
\item True or False: Electromagnetic radiation emitted from a nucleus is most likely in the form of gamma rays.
\item True or False: In metals, the conduction electrons form a degenerate Fermi gas, leading to a mean kinetic energy significantly higher than kT.
\item True or False: A helium-neon laser is the most effective choice for spectroscopy over a range of visible wavelengths.
\item True or False: The angular magnification of the refracting telescope is 6 when using an eye-piece lens with a focal length of 20 cm.
\item True or False: An observer moving at 0.50 c parallel to a 1.00 m rod will measure its length as 0.80 m due to relativistic effects.
\end{enumerate}

\section*{Long Form Response (20 marks)}
\textit{Answer each question with a clear and complete explanation.}
\begin{enumerate}[label=\arabic*.]
\setcounter{enumi}{10}
\item An airplane traveling at a constant speed of 100 m/s drops a payload. Analyze the motion of the payload 4.0 s after release, focusing on its velocity relative to the airplane.
\item Discuss the implications of a reversible thermodynamic process, focusing on how it affects the entropy of the system and its surroundings.
\end{enumerate}

\vfill
\noindent\textit{End of Paper}
\end{document}